\section{Pruebas funcionales y de integración}
\label{sec:pruebas}

\subsection{Principio rector}

\textbf{Las aserciones sobre dinero consultan PostgreSQL, no las respuestas
HTTP.} Un servicio que devolviera \codigo{201 Created} a todo sin escribir
nada pasaría cualquier prueba que solo mire códigos de estado.

Por la misma razón, las pruebas corren contra un PostgreSQL real y no contra
SQLite: lo que se verifica son garantías del motor (restricciones, triggers,
bloqueos y niveles de aislamiento) que otro motor no reproduce.

\subsection{Distribución}

\begin{table}[htbp]
\centering
\caption{Las 96 pruebas automatizadas del proyecto.}
\label{tab:pruebas}
\small
\begin{tabularx}{\textwidth}{@{}lrX@{}}
\toprule
\textbf{Fichero} & \textbf{N.º} & \textbf{Qué verifica} \\
\midrule
\codigo{test\_seguridad.py} & 19 &
  Prueba F, ordenada por riesgo de OWASP \\
\codigo{test\_invariantes\_esquema.py} & 17 &
  Ataca cada invariante por SQL directo, sin pasar por la aplicación \\
\codigo{test\_identidad.py} & 13 &
  Registro, autenticación, hashing y tokens \\
\codigo{test\_idempotencia.py} & 10 &
  Prueba C y el resto del protocolo, con un caso concurrente \\
\codigo{test\_transferencias.py} & 10 &
  Pruebas A y B, historial y conservación del saldo total \\
\codigo{test\_rate\_limit.py} & 10 &
  Cubo de fichas como unidad y aplicado a los endpoints \\
\codigo{test\_concurrencia.py} & 5 &
  Prueba D con las dos estrategias y ausencia de deadlocks \\
\codigo{test\_fallos.py} & 5 &
  Prueba E, con los tres fallos provocados de verdad \\
\codigo{test\_salud\_api.py} & 5 &
  Arranque, correlación y formato uniforme de error \\
\codigo{test\_migraciones.py} & 2 &
  Reversibilidad y ausencia de deriva entre modelos y esquema \\
\midrule
\textbf{Total} & \textbf{96} & \\
\bottomrule
\end{tabularx}
\end{table}

\subsection{Las seis pruebas exigidas}

\paragraph{A. Transferencia correcta.} Ana tiene S/\,100 y Luis S/\,20; Ana
envía S/\,30. Se comprueba en la base que Ana queda con 7\,000 céntimos, Luis
con 5\,000 y que existen exactamente dos asientos, $(-3000, +3000)$, ligados
por el mismo \codigo{transfer\_id}. Se verificó además con \codigo{curl}
contra el sistema levantado con Docker. \estado{verificado}

\paragraph{B. Fondos insuficientes.} Ana intenta enviar S/\,200 teniendo
S/\,100. La respuesta es 422, los saldos quedan \emph{idénticos} al céntimo y
la transferencia se registra con estado \codigo{rejected} y \textbf{cero
asientos}. Persistir el rechazo sin generar movimiento es intencionado: queda
en el historial sin tocar el dinero. \estado{verificado}

\paragraph{C. Solicitud duplicada.} Dos envíos con la misma clave de
idempotencia. La segunda respuesta es idéntica a la primera y lleva la
cabecera \codigo{Idempotent-Replay: true}; en el libro mayor hay un único par
de asientos nuevos. Se comprueba también que un rechazo por fondos es
igualmente idempotente y que el orden de los campos del JSON no altera la
huella. \estado{verificado}

\paragraph{D. Cien solicitudes concurrentes.} Detallada en la
sección~\ref{subsec:pruebaD}.

\paragraph{E. Fallos.} Detallada en la sección~\ref{subsec:pruebaE}.

\paragraph{F. Seguridad.} Diecinueve comprobaciones: operar sobre cuentas
ajenas, consultar operaciones ajenas, importes inválidos (cero, negativos,
decimales, fuera de rango, texto), campos extra, intentos de inyección SQL en
el alias de destino, y ausencia de traza interna en cualquier respuesta de
error. \estado{verificado}

\subsection{Prueba D, cien solicitudes concurrentes}
\label{subsec:pruebaD}

Una cuenta con S/\,100. Cien hilos, cada uno con su propia conexión, intentan
transferir S/\,30 a la vez desde una barrera común. El máximo posible de
aprobaciones es tres.

\begin{table}[htbp]
\centering
\caption{Prueba D en la suite automatizada. Tres ejecuciones consecutivas.}
\label{tab:pruebaD}
\small
\begin{tabular}{@{}llrrrr@{}}
\toprule
\textbf{Ejec.} & \textbf{Estrategia} & \textbf{Aprobadas} &
\textbf{Rechazadas} & \textbf{Reint. agotados} & \textbf{Saldo final} \\
\midrule
1 & \codigo{pessimistic}   & 3 & 97 & 0 & 1\,000 ctms \\
1 & \codigo{serializable}  & 3 & 94 & 3 & 1\,000 ctms \\
2 & \codigo{pessimistic}   & 3 & 97 & 0 & 1\,000 ctms \\
2 & \codigo{serializable}  & 3 & 92 & 5 & 1\,000 ctms \\
3 & \codigo{pessimistic}   & 3 & 97 & 0 & 1\,000 ctms \\
3 & \codigo{serializable}  & 3 & 93 & 4 & 1\,000 ctms \\
\bottomrule
\end{tabular}
\end{table}

En las seis ejecuciones se aprueban exactamente tres transferencias, el saldo
final es 1\,000 céntimos, ningún saldo queda negativo y la reconciliación
cuadra. Las dos estrategias son correctas; se diferencian en que
\codigo{SERIALIZABLE} devuelve un 503 a entre el 3 y el 5\,\% de las
peticiones, que no es una respuesta sobre su dinero sino sobre el estado del
sistema.

Esta prueba es la que destapó la actualización perdida descrita en la
sección~\ref{subsec:lostupdate}.

\subsection{Prueba E, fallos}
\label{subsec:pruebaE}

Los tres fallos se provocan de verdad, no se simulan con comentarios.

\paragraph{Antes del \codigo{COMMIT}.} Se inyecta una excepción después de
mover los saldos y antes de confirmar. Resultado verificado en la base:
saldos idénticos, cero asientos nuevos y \textbf{ninguna} clave de
idempotencia registrada. Como la clave tampoco se persistió, el reintento del
cliente se procesa como una petición nueva y la transferencia se ejecuta una
sola vez.

\paragraph{Después del \codigo{COMMIT}, antes de responder.} Se rompe la
construcción de la respuesta una vez confirmada la transacción. El dinero
\emph{sí} se movió, el cliente recibe un 500 y no sabe si pagó. Al reintentar
con la misma clave recibe 201 con \codigo{Idempotent-Replay: true} y el estado
de la base no cambia ni un céntimo. Es exactamente el escenario para el que
existe el diseño de la sección~\ref{sec:idempotencia}.

Hay una contraprueba: dos peticiones idénticas con claves \emph{distintas} sí
mueven el dinero dos veces. Deja claro que lo que evita el duplicado es la
clave y no ninguna otra propiedad del sistema.

\paragraph{Caída del proceso de base de datos.} No se simula. Se obtiene el
identificador del proceso servidor de la transacción en curso y se mata desde
otra conexión con \codigo{pg\_terminate\_backend}, con los asientos ya
escritos y sin confirmar. Es lo más parecido a un corte de luz que se puede
provocar de forma reproducible. El \codigo{COMMIT} falla, el estado de la base
queda idéntico al previo y la reconciliación sigue cuadrando. Una segunda
prueba comprueba que el servicio queda operativo: la verificación previa del
pool descarta la conexión muerta y la siguiente transferencia se completa con
normalidad. \estado{verificado}

\subsection{Integración continua}

El flujo de trabajo levanta PostgreSQL como servicio del runner, por lo que
las pruebas de integración y concurrencia se ejecutan de verdad en CI y no
solo las unitarias. Un segundo trabajo comprueba los tipos y la compilación
del frontend. El CI \textbf{no} se usa para medir rendimiento: los runners
compartidos no son un entorno fiable para eso.
